\section{Nonzero transfer and scalar reassembly}
\label{sec:tpc27-nonzero-scalar}

The new work at additive frequency zero is complete.  We now transfer
the inherited exact-conductor theorem to the same physical packet and
then restore the three calibration channels in
\eqref{eq:tpc27-ZZ-EZ-ZE-EE}.

\subsection{The inherited dynamic conductor region}

Write
\begin{equation}
 Q=X^{\beta+o(1)},
 \qquad J=X^{1-\beta+o(1)},
 \qquad S=X^{s+o(1)}.
 \label{eq:tpc27-base-exponents}
\end{equation}
For the selected square \(u,v\le S\), the common exact-conductor
ledger of TPC-22 at \(S=R\) and TPC-23 at \(S>R\) is
\begin{equation}
 \eta_S(\beta,s)
 =\min\left\{
 1-\frac{3\beta}{2},
 \frac{\beta+1-4s}{2},
 \frac{3-\beta-6s}{4}
 \right\}.
 \label{eq:tpc27-etaS}
\end{equation}
The corresponding nonzero theorem assumes
\begin{equation}
 s<1-\beta,
 \qquad \eta_S(\beta,s)>0
 \label{eq:tpc27-dynamic-region}
\end{equation}
with fixed margins, together with the primitive row interface,
divisor-independent off-diagonal mask, no-wrap condition, and the
stated smooth Fourier-tail control
\citep{WangTPC23}.  The first inequality is the strict Poisson-short
condition
\begin{equation}
 HS\le JX^{-\eta_0}
 \label{eq:tpc27-Poisson-short}
\end{equation}
for some fixed \(\eta_0>0\), after adjusting fixed constants.

At the boundary \(S=R\) we use the static theorem of TPC-22; in the
genuinely dynamic case \(S>R\) we use TPC-23.  The theorem acts on the
retained nonzero joint packet.  Its global
principal multiplicative-character cell is bounded absolutely by
\(Q^2X^\eps\), while every nonprincipal exact-conductor cell has a
fixed-power saving from \eqref{eq:tpc27-etaS}.  It does not include the
additive zero kernel; that is why
\cref{thm:tpc27-physical-base-zero} was needed.

\subsection{Transfer of the nonzero spectrum}

Let \(\cN_S^{ZZ,\ne0}\) be the complete physical contribution of all
nonzero additive frequencies from the centered product
\(Z_{m_\alpha,S}Z_{m_\gamma,S}\), after the additive zero term
\eqref{eq:tpc27-physical-ZZ-zero} has been set aside.

\begin{lemma}[Physical transfer to the dynamic interface]
\label{lem:tpc27-physical-nonzero-transfer}
After splitting the \(O_h(1)\) residue cells, the physical nonzero
packet has a continuous representation
\begin{equation}
 \gamma_\alpha^{(1)}\gamma_\gamma^{(2)}
 \cW_{\alpha,\gamma}(j/J)
 =\int
 x_{\boldsymbol\tau}(\alpha)
 y_{\boldsymbol\tau}(\gamma)
 \Phi_{\boldsymbol\tau}(j/J)
 \,d\nu(\boldsymbol\tau),
 \label{eq:tpc27-continuous-transfer}
\end{equation}
such that, for every fixed integer \(M\ge0\),
\begin{align}
 \int(1+|\boldsymbol\tau|)^M
 \|x_{\boldsymbol\tau}\|_1
 \|y_{\boldsymbol\tau}\|_1
 \,d|\nu|(\boldsymbol\tau)
 &\ll_{M,\mathscr D}Q^2,
 \label{eq:tpc27-transfer-l1}\\
 \int(1+|\boldsymbol\tau|)^M
 \|x_{\boldsymbol\tau}\|_2
 \|y_{\boldsymbol\tau}\|_2
 \,d|\nu|(\boldsymbol\tau)
 &\ll_{M,\mathscr D}Q\log(2L).
 \label{eq:tpc27-transfer-l2}
\end{align}
Every component has the same divisor-independent mask, row scales,
and no-wrap margins.  The applicable TPC-22 static or TPC-23 dynamic
mean, discrepancy, and Fourier-tail bounds may be integrated
componentwise.
For a fixed cell \(j\equiv\varrho\pmod{H_0}\), the bounded factor
\(\Xi_{\alpha,\gamma}(\varrho)\) is absorbed into the row-pair mask.
It remains independent of every divisor, additive frequency, and
conductor variable.
\end{lemma}

\begin{proof}
This is the continuous log-Mellin separation proved for the same
physical provenance in TPC-25.  Each fixed smooth row--orbit factor has
a Schwartz Mellin transform.  Its row phase and orbit phase have
modulus one, so the row norms are exactly
\cref{eq:tpc27-row-l1,eq:tpc27-row-l2} up to the Schwartz density.
Schwartz moments prove
\cref{eq:tpc27-transfer-l1,eq:tpc27-transfer-l2}.  Smooth separation
does not touch the generic mask.

The constants in the applicable TPC-22 or TPC-23 component estimates
grow at most
polynomially in finitely many seminorms of
\(\Phi_{\boldsymbol\tau}\), while the weighted moments above are finite.
Apply the triangle inequality and Tonelli's theorem to finite Riemann
sums, then use dominated convergence.  The physical Fourier tail is
automatic for this provenance by the rapid-decay argument of TPC-26;
the base support \(u,v\le S\) is a subfamily of the polynomial divisor
box treated there.  This proves the transfer.
\end{proof}

\begin{theorem}[Physical nonzero centered-base closure]
\label{thm:tpc27-physical-nonzero}
Assume \eqref{eq:tpc27-dynamic-region} with fixed margins and all
inherited primitive/no-wrap interfaces.  Then for every fixed \(B>0\),
\begin{equation}
 \boxed{
 \cN_S^{ZZ,\ne0}
 \ll_{B,\mathscr D}XQ(\log X)^{-B}.}
 \label{eq:tpc27-physical-nonzero}
\end{equation}
\end{theorem}

\begin{proof}
Apply \cref{lem:tpc27-physical-nonzero-transfer}.  On every Mellin
component, TPC-20's two-scale completion removes the single-indicator
and constant nonzero spectra under
\eqref{eq:tpc27-Poisson-short}.  If \(S=R\), apply the TPC-22 static
exact-conductor theorem; if \(S>R\), apply the TPC-23 dynamic theorem.
Their common ledger \eqref{eq:tpc27-etaS} gives the same fixed-power
saving.  The relevant Ramanujan-mean theorem bounds the global
principal character cell by \(Q^2X^\eps\), a fixed-power saving from
\(Q^2J\) because \(J\) is a fixed positive power.  Integrate the
component estimates using \eqref{eq:tpc27-transfer-l2}, the mean using
\eqref{eq:tpc27-transfer-l1}, and choose the Fourier-tail decay order
after \(B\).  Every contribution is then
\(O_B(XQ(\log X)^{-B})\).
\end{proof}

\subsection{The three calibration channels}

The fact that \(E_S(m)\) depends on the row prevents it from being
removed as one global scalar.  Uniform pointwise calibration and the
actual physical \(\ell^1\)-mass nevertheless suffice.

\begin{lemma}[Orbit mass of one centered prefix]
\label{lem:tpc27-centered-orbit-mass}
For every retained row and every uniformly bounded physical orbit
test supported on an interval of length \(O(J)\),
\begin{equation}
 \sum_{j\in\Z}
 \left|F(j/J)Z_{m,S}(j)\right|
 \ll_{\mathscr D}(J+S)(\log(2X))^4.
 \label{eq:tpc27-centered-orbit-mass}
\end{equation}
The same bound holds after imposing \((j,H)=1\) and a fixed residue
cell modulo \(H_0\).
\end{lemma}

\begin{proof}
For \((u,mH)=1\), the congruence \(u\mid mj+h\) selects one class
modulo \(u\).  Hence its intersection with an interval of length
\(O(J)\) has \(O(J/u+1)\) points.  From
\eqref{eq:tpc27-Zms},
\begin{align*}
 \sum_j|F(j/J)Z_{m,S}(j)|
 &\ll
 \sum_{u\le S}|b_R(u)|\left(\frac Ju+1\right).
\end{align*}
TPC-24's prefix mass bound gives
\(\sum_{u\le S}|b_R(u)|/u\ll(\log(2X))^3\).  Pointwise, the defining
formula \eqref{eq:tpc27-lambda-prime} gives
\[
 |\lambda'_R(u)|
 \le \frac{u}{\varphi(u)}
 \sum_{b\le R/u}\frac{\mu^2(b)}{\varphi(b)}
 \ll (\log(2X))^2
\]
on its squarefree support; also \(|a(u)|\le\log(2X)\).  Hence
\(\sum_{u\le S}|b_R(u)|\ll S(\log(2X))^2\).  Combining these
estimates proves \eqref{eq:tpc27-centered-orbit-mass}.  Additional
fixed congruence restrictions only reduce the number of orbit points.
\end{proof}

Let \(\cE_S^{EZ}\), \(\cE_S^{ZE}\), and \(\cE_S^{EE}\) denote the
three physical terms obtained from the last three summands of
\eqref{eq:tpc27-ZZ-EZ-ZE-EE}.

\begin{theorem}[Scalar calibration closure]
\label{thm:tpc27-scalar-closure}
Assume the actual physical provenance and the strict Poisson-short
condition \eqref{eq:tpc27-Poisson-short}.  Then, for every fixed
\(B>0\),
\begin{equation}
 \boxed{
 |\cE_S^{EZ}|+|\cE_S^{ZE}|+|\cE_S^{EE}|
 \ll_{B,\mathscr D}XQ(\log X)^{-B}.}
 \label{eq:tpc27-scalar-closure}
\end{equation}
\end{theorem}

\begin{proof}
Choose the exponent \(A\) in \eqref{eq:tpc27-E-bound} after the
requested \(B\).  Delete the bounded mask, residue factor, and orbit
test in absolute value.  By
\cref{eq:tpc27-row-l1,eq:tpc27-centered-orbit-mass},
\begin{align}
 |\cE_S^{EZ}|+|\cE_S^{ZE}|
 &\ll_A
 Q^2(J+S)(\log(2X))^4(\log X)^{-A},
 \label{eq:tpc27-EZ-bound}\\
 |\cE_S^{EE}|
 &\ll_A
 JQ^2(\log X)^{-2A}.
 \label{eq:tpc27-EE-bound}
\end{align}
Condition \eqref{eq:tpc27-Poisson-short} gives \(J+S\asymp J\), and
\(JQ^2\asymp XQ\).  Increase \(A\) after the requested \(B\) to
obtain \eqref{eq:tpc27-scalar-closure}.
\end{proof}

This proof treats \(EZ,ZE,EE\) as genuine row-dependent channels.  It
does not replace them by a single scalar or infer their smallness from
the nonzero conductor theorem.
